\section{A Planning Stack}

For teaching purposes, it is useful to separate planning into layers:
\begin{enumerate}[nosep]
  \item \textbf{Outcome layer:} what end state is desired?
  \item \textbf{Session layer:} what concrete session will move that outcome?
  \item \textbf{Entry layer:} what exact first action starts the session?
  \item \textbf{Environment layer:} what physical arrangement lowers entry
        friction?
  \item \textbf{Trigger layer:} when and where is the decision window expected?
\end{enumerate}

\begin{proposition}[Execution-choice compression]
If a plan resolves $k \geq 1$ future micro-decisions before the decision
window, then the execution choice complexity at the window is reduced from $D$
to at most $D-k$.
\end{proposition}

\begin{proof}
By construction, each resolved micro-decision is one less action-relevant item
to be decided at execution time. Therefore resolving $k$ such items before the
window reduces the remaining decision count by at least $k$, so the future
decision complexity is at most $D-k$.
\end{proof}

\begin{corollary}[Binary execution advantage, qualified]
If a plan reduces execution complexity to a binary choice---execute the
predefined action or refuse it---then the in-the-moment execution burden is
weakly lower than under any otherwise similar plan with strictly larger
execution complexity.
\end{corollary}

\begin{proof}
A binary choice has execution complexity $D=1$. Any otherwise similar plan with
strictly larger complexity has $D>1$. By the decision-load assumption, greater
execution complexity weakly increases in-the-moment burden. Therefore the
binary-choice plan weakly lowers that burden.
\end{proof}

\begin{discussion}
This result is narrower and more honest than the older formulation. It does not
claim that binary framing guarantees compliance. It claims only that, within
the model, reducing the number of unresolved real-time decisions reduces the
decision burden that must be paid at the window.
\end{discussion}
